% !TeX program = pdflatex
\documentclass[11pt]{amsart}

\usepackage[dvipsnames]{xcolor}
\usepackage{enumitem}
\usepackage{amssymb,graphicx,mathtools,microtype}
\usepackage{tikz}
\usetikzlibrary{calc,arrows.meta}
\usepackage[margin=1in]{geometry}
\geometry{letterpaper}
\usepackage{etoolbox,needspace}
\usepackage[square,longnamesfirst]{natbib}

\theoremstyle{plain}
\newtheorem{theorem}{Theorem}[section]
\newtheorem{proposition}[theorem]{Proposition}
\newtheorem{corollary}[theorem]{Corollary}
\newtheorem{lemma}[theorem]{Lemma}
\theoremstyle{definition}
\newtheorem{problem}{Problem}

\newcommand{\Z}{\mathbb Z}
\newcommand{\R}{\mathbb R}
\renewcommand{\P}{\mathbb P}
\newcommand{\E}{\mathbb E}
\newcommand{\eps}{\varepsilon}
\newcommand{\ind}{\mathbf 1}
\newcommand{\dd}{\mathrm d}
\newcommand{\defn}[1]{\textbf{#1}}
\setcounter{secnumdepth}{2}

\makeatletter
\patchcmd{\@settitle}{\uppercasenonmath\@title}{\Large}{}{}
\patchcmd{\@setauthors}{\MakeUppercase}{\large}{}{}
\makeatother

\usepackage[citecolor=PineGreen,colorlinks=true,
linkcolor=RoyalBlue,urlcolor=BrickRed]{hyperref}
\hypersetup{pdftitle={The limit shape of centrally excited random walk},
pdfauthor={Ahmed Bou-Rabee and Yuval Peres}}

\title{The limit shape of centrally excited random walk}
\author{Ahmed Bou-Rabee}
\address{Department of Mathematics, University of Pennsylvania,
Philadelphia, PA 19104, USA}
\email{ahmedmb@sas.upenn.edu}
\author{Yuval Peres}
\address{Beijing Institute of Mathematical Sciences and Applications,
Beijing 101408, China}
\email{yperes@bimsa.cn}
\subjclass[2020]{Primary 60K35; Secondary 60G50, 60G42, 31C20}
\keywords{Centrally excited random walk, limit shape, local time,
recurrence, fluctuations}
\date{}

\begin{document}

\begin{abstract}
A centrally excited random walk takes a step with a mean of fixed magnitude toward
the origin when it first leaves a nonzero site, and takes simple random walk
steps on later visits. We prove that, in every dimension at least two,
the region visited in the first $n$ steps, rescaled by $n^{1/(d+1)}$,
converges almost surely to a Euclidean ball. We determine its radius
and show that the rescaled local times converge uniformly to a cone.
We also give polynomial upper bounds on the inner and outer fluctuations.
The proof represents local times by an inward-drift potential and uses
positivity at an unvisited point nearest the origin to compare the range with a ball.
\end{abstract}

\maketitle

\begin{figure}[!b]
\centering
\includegraphics[width=0.76\textwidth]{Figures/first_departure.pdf}
\caption{Conditional means of the next departure at each site,
before the walk (left) and after $120$ steps (right), with $\eps=0.4$.
Each blue arrow disappears when the walk first leaves its site.
Gray edges trace the walk from the origin (open circle) to its
current position (red dot).}
\label{fig:mechanism}
\end{figure}

\section{Introduction}\label{sec:introduction}

A \defn{centrally excited random walk} starts at the origin of the square
lattice. When it first leaves a site other than the origin, its step has
mean $-\eps x/|x|$, where $x$ is the site and $0<\eps<1/2$ is fixed.
On later visits it chooses each neighbor with equal probability
(Figure~\ref{fig:mechanism}).
We prove that the visited region has a circular limit shape, with radius
asymptotic to $(3n/(4\pi\eps))^{1/3}$ after $n$ steps. The local time
decreases linearly with the distance from the origin, to first order.

Kozma posed the ball-shape problem in the 2007 Oberwolfach problem
session \citep[p.~1552]{Oberwolfach2007}.
\citet*{BenjaminiDuminilCopinKozmaLucas2020} later described the
conjectured radius $n^{1/(d+1)}$ in their introduction.
We establish the ball shape for a constant Euclidean inward mean in
every dimension $d\geq2$, and give quantitative bounds on its fluctuations.
Figure~\ref{fig:ranges} shows the growth in dimensions two and three.



\subsection{Main results}

Let $d\geq2$ be an integer, let $0<\eps<1/d$, and let
$e_1,\ldots,e_d$ be the coordinate vectors of $\Z^d$.
We start the walk at $X_0=0$. On the first departure from $x\ne0$,
the transition probabilities are
\begin{equation}\label{eq:kernel}
q_x(\pm e_i)=\frac1{2d}\mp\frac\eps2\frac{x_i}{|x|}
\qquad (1\leq i\leq d)\, .
\end{equation}
Every later departure, and every departure from the origin, is a simple
symmetric random walk step. These conditional transition probabilities
specify the law of the process given its entire past.



The \defn{departure local time} and the two ranges are
\begin{equation}\label{eq:local-range}
\ell_n(x)=\sum_{j=0}^{n-1}\ind_{\{X_j=x\}} ,
\qquad A_n=\{x:\ell_n(x)>0\} ,
\qquad V_n=\{X_0,\ldots,X_n\}\, .
\end{equation}
Thus $V_n=A_n\cup\{X_n\}$ and $|V_n\setminus A_n|\leq1$.
We write $B(0,r)$ for the open Euclidean ball in $\R^d$ and
$\omega_d$ for the volume of the unit ball. Let
\begin{equation}\label{eq:scale}
N=n^{1/(d+1)} ,\qquad c_0=2d\eps ,\qquad
a=\left(\frac{d+1}{2d\eps\omega_d}\right)^{1/(d+1)}\, .
\end{equation}

\Needspace{18\baselineskip}
\begin{theorem}[Ball shape]\label{thm:shape}
Let $d\geq2$ be an integer, let $0<\eps<1/d$, and let $(X_n)_{n\geq0}$
be centrally excited random walk on $\Z^d$, started at the origin,
with first-departure probabilities
$q_x(\pm e_i)=1/(2d)\mp\eps x_i/(2|x|)$ for $x\ne0$ and simple
random walk steps on every other departure. Let
$a=((d+1)/(2d\eps\omega_d))^{1/(d+1)}$ and $N=n^{1/(d+1)}$.
Then the following statements hold almost surely.
\begin{enumerate}[label=\textup{(\roman*)}]
\item For every $0<\eta<a$, for all sufficiently large integers $n$,
\begin{equation}\label{eq:shape}
B(0,(a-\eta)N)\cap\Z^d\subseteq A_n\subseteq V_n
\subseteq B(0,(a+\eta)N)\cap\Z^d\, .
\end{equation}
\item The departure local times satisfy
\begin{equation}\label{eq:profile-limit}
\sup_{x\in\Z^d}\left|
\frac{\ell_n(x)}N-2d\eps\left(a-\frac{|x|}N\right)_+
\right|\longrightarrow0\, .
\end{equation}
\item The range satisfies $|V_n|/N^d\longrightarrow\omega_da^d$.
For every fixed $x\in\Z^d$, the local time satisfies
$\ell_n(x)/N\longrightarrow2d\eps a$; in particular, every vertex
is visited infinitely often.
\end{enumerate}
\end{theorem}

\begin{figure}[t]
\centering
\includegraphics[width=0.76\textwidth]{Figures/cerw_growth.pdf}
\caption{Departure ranges in dimension two (left) and central slices
$x_3=0$ of the ranges in dimension three (right). The top row shows
$10^8$ steps; the bottom row shows $2500$ steps of the same trajectories.
Blue to red records earlier to later first departures.
The drift strengths are $\eps=0.4$ and $\eps=0.25$, respectively.
Coordinates are divided by $a n^{1/(d+1)}$, with a common window in each row.}
\label{fig:ranges}
\end{figure}

Our second theorem bounds the distance from the ball.
To compare volumes, replace each site $x\in A_n$ by the half-open
unit cell $C_x=x+[-1/2,1/2)^d$, and let
$D_n=\bigcup_{x\in A_n}C_x$. For $L=\log(n+2)$, write
\begin{equation}\label{eq:rates}
Q=\begin{cases}
(L/N)^{1/2},&d=2,\\
(L/N)^{d/(2d-1)},&d\geq3
\end{cases}\, .
\end{equation}

\Needspace{22\baselineskip}
\begin{theorem}[Fluctuation bounds]\label{thm:fluctuations}
Let $d\geq2$ be an integer, let $0<\eps<1/d$, and let $(X_n)_{n\geq0}$
be centrally excited random walk on $\Z^d$, started at the origin,
with first-departure probabilities
$q_x(\pm e_i)=1/(2d)\mp\eps x_i/(2|x|)$ for $x\ne0$ and simple
random walk steps on every other departure. Let
$a=((d+1)/(2d\eps\omega_d))^{1/(d+1)}$, let $N=n^{1/(d+1)}$, and let
$Q=(\log(n+2)/N)^{1/2}$ for $d=2$ and
$Q=(\log(n+2)/N)^{d/(2d-1)}$ for $d\geq3$.
For every $p>0$, there are constants $C,n_0>0$, depending only on
$d,\eps,p$, such that for every integer $n\geq n_0$, with probability
at least $1-Cn^{-p}$,
\begin{equation}\label{eq:sandwich}
B(0,(a-CQ)N)\cap\Z^d\subseteq A_n\subseteq V_n
\subseteq B(0,(a+CQ^{1/d}L)N)\cap\Z^d\, ,
\end{equation}
where $L=\log(n+2)$. On the same event,
\begin{align}
|(D_n/N)\mathbin{\triangle}B(0,a)|
+\left|\frac{|A_n|}{N^d}-\omega_da^d\right|&\leq CQ\, ,
\label{eq:volume}\\
\sup_{x\in\Z^d}\left|
\frac{\ell_n(x)}N-2d\eps(a-|x|/N)_+\right|&\leq CQ^{1/d}\, .
\label{eq:profile-rate}
\end{align}
All these bounds hold eventually almost surely, with deterministic
constants and a finite random starting time.
\end{theorem}

On the event in Theorem~\ref{thm:fluctuations}, the Euclidean Hausdorff
distance satisfies
\begin{equation}\label{eq:hausdorff}
d_H(N^{-1}V_n,\overline B(0,a))\leq C
\begin{cases}
n^{-1/12}L^{5/4},&d=2,\\
n^{-1/((d+1)(2d-1))}L^{2d/(2d-1)},&d\geq3
\end{cases}\, .
\end{equation}
The inner-radius error in the plane is at most
$Cn^{1/6}\sqrt{\log(n+2)}$ in lattice units.
The bounds above are upper bounds; the orders of the fluctuations
remain open.

\begin{figure}[t]
\centering
\includegraphics[width=0.82\textwidth]{Figures/cerw_radial_profiles.pdf}
\caption{Local times after $10^8$ steps, normalized by $2d\eps r_n$,
where $r_n=a n^{1/(d+1)}$. Above: the planar field and the
three-dimensional slice $x_3=0$; purple to yellow indicates increasing
local time on a common scale. Below: shell averages (blue) and
$(1-s)_+$ (dashed black), where $s=|x|/r_n$.
Shells have width $0.025r_n$ and include all lattice sites in the full
dimension, including unvisited sites. Parameters are as in Figure~\ref{fig:ranges}.}
\label{fig:profiles}
\end{figure}

\subsection{Proof overview}

The proof begins with a potential representation of local times.
For a bounded measurable set $D\subset\R^d$, let
\begin{equation}\label{eq:potential-intro}
U_D(y)=\frac{2\eps}{\omega_d}\int_D
\frac{v}{|v|}\cdot\frac{v-y}{|v-y|^d}\dd v\, .
\end{equation}
The value of $v/|v|$ at zero does not affect this integral.
The potential $U_{D_n}$ approximates $\ell_n$ uniformly.
Its spherical average at radius $r$ is proportional to the weighted
volume of $D_n$ outside that radius.

In Section~\ref{sec:coarse}, we use this identity to prove that the
range has radius of order $N$. We truncate the exact lattice potential
kernel at a radial level. The resulting test is harmonic outside a
shell of bounded width, in the plane as well as in higher dimensions.
The test controls the weighted exterior volume. An occupation estimate
on intervals then excludes long excursions through a small set of sites.

The shape is determined in Section~\ref{sec:contact}.
Let $b$ be the radius of the largest centered ball contained in $D_n$,
and let $y_0$ be a boundary point with $|y_0|=b$.
The ball potential is $U_{B(0,b)}(y)=2d\eps(b-|y|)_+$.
Every part of $D_n$ outside this ball contributes positively to
$U_{D_n}(y_0)$. The local time at a neighboring unvisited site is zero,
so concentration bounds this contribution. A local variance estimate
improves the resulting bound on the excess volume to $CN^dQ$.
The martingale for $|X_n|^2$ then determines $b$ to within $CNQ$ of $aN$.
In Section~\ref{sec:outer}, we use the exterior-volume estimate again
to obtain the outer fluctuation bound. Theorem~\ref{thm:shape} follows
from Theorem~\ref{thm:fluctuations}.

\subsection{Related work}

The model is also called \emph{origin-excited random walk}
\citep[Introduction]{DemboGroismanHuangSidoravicius2021}.
\citet[Section~2]{Kozma2013} discusses excitation toward the center
and reports recurrence in every dimension, citing his unpublished
work from 2006. The name \emph{centrally excited random walk} occurs
in the title of that manuscript, recorded as reference~[K2] by
\citet{DemboHuangSidoravicius2014}.
The inward excitation differs from the fixed directional bias in the
excited random walk of \citet{BenjaminiWilson2003}.

\citet*[Theorems~1.4 and~1.9]{DemboGroismanHuangSidoravicius2021}
prove an averaging principle and a shape theorem for a growth process
in Euclidean space with memory. Their Figure~2 shows lattice walks
with several inward-excitation rules.
\citet{DemboYang2025Flow} obtain a
flow scaling limit for growth driven by reflected Brownian motion, and
\citet{DemboYang2025KPZ} derive a KPZ-type fluctuation equation for
smoothed growth driven by reflected Brownian motion.

Ball shapes also occur in internal diffusion-limited aggregation
\citep{LawlerBramsonGriffeath1992} and rotor-router aggregation
\citep{LevinePeres2009}; see
\citet{LevinePeres2017} for a survey.
For a single planar Eulerian walker, \citet{KapriDhar2009}
conjectured a circular visited region from simulations.
\citet[Theorem~1.1]{BouRabeePeres2026} prove a deterministic convex
limit shape for clockwise rotors with independent uniform initial
directions. The identification of that shape with a disk remains open.
The potential here comes from the explicit first-departure mean,
so its radial form is specific to centrally excited random walk.

The constant Euclidean magnitude in our hypothesis matters.
A mean proportional to $-x/\|x\|_1$ has an angularly varying Euclidean
magnitude and is not covered by the ball conclusion.
Theorems~\ref{thm:shape} and~\ref{thm:fluctuations} concern $d\geq2$.
Their concentration argument does not extend to dimension one, where
the martingale displacement remains of the same order as the range.

\subsection{Notation and conventions}

We use $|x|$ for Euclidean norm, $|A|$ for cardinality of a lattice set,
and $|D|$ for Lebesgue measure of a subset of $\R^d$.
We write $x\sim y$ when $x,y\in\Z^d$ are nearest neighbors, and
$t_+=\max(t,0)$ for $t\in\R$.
Throughout the proof, $d\geq2$ and $\eps\in(0,1/d)$ are fixed.
For each probability exponent $p>0$, constants $c,C>0$ depend only on
$d,\eps,p$ unless further dependencies are displayed. Constants may
change between estimates. All estimates concern integers $n\geq2$;
statements requiring large $n$ have deterministic thresholds.
The symbols $N,L,Q,a,c_0$ have the meanings given above.

We use the additional occupation quantities
\begin{equation}\label{eq:occupation}
R_n=|A_n| ,\qquad M_n=\max_x\ell_n(x) ,\qquad
H_n=\max_{0\leq j\leq n}|X_j|\, .
\end{equation}
Let $\widetilde\ell_n$ equal $\ell_n(x)$ on $C_x$, including the
value zero on unoccupied cells. Let $u_x=x/|x|$ for $x\ne0$ and
$u_0=0$, for lattice or real $x$. The indicator of a first departure
at time $j$ is $I_j=\ind_{\{X_j\notin A_j\}}$.
The filtration is $\mathcal F_j=\sigma(X_0,\ldots,X_j)$.
For a martingale $Z$, the notation $\langle Z\rangle_n$ denotes
its predictable quadratic variation through time $n$.

\section{The local-time potential}\label{sec:local}
We represent occupation times by a potential whose spherical averages
measure the part of the range outside a ball.

Let $P$ be the simple random walk transition operator on $\Z^d$.
For $d=2$, let $b$ be the potential kernel normalized by $b(0)=0$.
For $d\geq3$, let $b=-G$, where
$G(x)=\sum_{j\geq0}P^j(0,x)$ is the Green function.
In both cases $(P-I)b=\ind_{\{0\}}$.
The lattice potential estimates of
\citet[Theorem~4.3.1, Corollary~4.3.3, and Theorem~4.4.4]{LawlerLimic2010}
give, as $|x|\to\infty$,
\begin{equation}\label{eq:kernel-asymptotics}
b(x)=\begin{cases}
\dfrac2\pi\log|x|+\kappa+O(|x|^{-2}),&d=2,\\[3pt]
-\dfrac{2}{(d-2)\omega_d}|x|^{2-d}+O_d(|x|^{-d}),&d\geq3
\end{cases}\, ,
\end{equation}
where $\kappa$ is the planar potential-kernel constant.
For the central difference $D_i b(x)=(b(x+e_i)-b(x-e_i))/2$,
\begin{equation}\label{eq:gradient}
Db(x)=\frac2{\omega_d}\frac{x}{|x|^d}+O_d(|x|^{-d}) ,
\qquad |b(x+e)-b(x)|\leq C_d(1+|x|)^{1-d}\, .
\end{equation}
The first estimate concerns $x\ne0$; the second concerns every lattice
site and every unit coordinate vector with either sign.

We use the following consequence of \citet[Theorem~1.6]{Freedman1975}.
If the increments of a real martingale $Z$ have absolute value at most
$B$, then, for deterministic $v,t>0$,
\begin{equation}\label{eq:freedman}
\P\bigl(|Z_n-Z_0|\geq t,\ \langle Z\rangle_n\leq v\bigr)
\leq2\exp\left(-\frac{t^2}{2(v+Bt)}\right)\, .
\end{equation}
We apply this inequality at dyadic variance bounds and take polynomial
union bounds over deterministic targets and time intervals.
The resulting estimates can then be used at points and intervals
selected from the path.

For $m\geq0$, define the error scale
\begin{equation}\label{eq:error-scale}
e_n(m)=\begin{cases}
\sqrt m L+L,&d=2,\\
\sqrt{mL}+L,&d\geq3
\end{cases}
\qquad
\lambda_n=\begin{cases}L^2,&d=2,\\L,&d\geq3
\end{cases}\, .
\end{equation}
For integers $0\leq s<t\leq n$, let
\begin{equation}\label{eq:interval-def}
\ell_{s,t}(x)=\sum_{s\leq j<t}\ind_{\{X_j=x\}} ,\qquad
M_{s,t}=\max_x\ell_{s,t}(x) ,\qquad
k_{s,t}=\sum_{s\leq j<t}I_j\, .
\end{equation}

\begin{lemma}\label{lem:local}
Let $d\geq2$ be an integer, let $0<\eps<1/d$, and let $(X_j)_{j\geq0}$ be
centrally excited random walk started at the origin, with first-departure mean $-\eps u_x$
and probabilities $q_x(\pm e_i)=1/(2d)\mp\eps u_x^i/2$,
with simple symmetric steps on later departures and at the origin.
For every $p>0$ there is $C=C(d,\eps,p)$ such that, for every integer
$n\geq2$, with probability at least $1-Cn^{-p}$, the following estimates
hold simultaneously:
\begin{align}
M_n&\leq C_d\eps R_n^{1/d}+C\lambda_n\, ,\label{eq:M}\\
M_{s,t}&\leq C_d\eps k_{s,t}^{1/d}+C\lambda_n
\qquad(0\leq s<t\leq n)\, ,\label{eq:interval}\\
|\widetilde\ell_n(y)-U_{D_n}(y)|&\leq C e_n(M_n)
\qquad(y\in\R^d,\ |y|\leq2n)\, .\label{eq:approx}
\end{align}
Here $U_D(y)=(2\eps/\omega_d)\int_D u_v\cdot(v-y)|v-y|^{-d}\dd v$,
$L=\log(n+2)$, and $\lambda_n,e_n$ have the values displayed above.
\end{lemma}

\begin{proof}
Dynkin's formula for $b(X_j-y)$ gives a martingale $\mathcal M^y$ with
\begin{equation}\label{eq:dynkin}
\ell_n(y)=b(X_n-y)-b(-y)
+\eps\sum_{x\in A_n}u_x\cdot Db(x-y)-\mathcal M_n^y\, .
\end{equation}
The endpoint difference is $O(L)$ for $d=2$ and $O_d(1)$ for $d\geq3$,
uniformly for lattice targets $|y|\leq3n$.
The increments of $\mathcal M^y$ are bounded, and
\begin{equation}\label{eq:bracket}
\langle\mathcal M^y\rangle_n\leq C_d B_y ,\qquad
B_y=\sum_x\ell_n(x)(1+|x-y|)^{2-2d}\, .
\end{equation}
The sum of the weights over $|x-y|\leq4n$ is $O(L)$ for $d=2$
and $O_d(1)$ for $d\geq3$. Applying \eqref{eq:freedman} at dyadic
bounds for $M_{s,t}\leq n$, and taking a union bound over the
$O_d(n^{d+2})$ interval--target pairs, gives
\begin{equation}\label{eq:interval-mart}
|\mathcal M_t^y-\mathcal M_s^y|\leq C e_n(M_{s,t})\, .
\end{equation}
For every set $E$ of $k$ lattice sites, filling concentric shells gives
\begin{equation}\label{eq:packing-lattice}
\sup_y\sum_{x\in E}(1+|x-y|)^{1-d}\leq C_d k^{1/d}\, .
\end{equation}
Subtract \eqref{eq:dynkin} at $s$ and $t$ and maximize over $y$.
The source sum contains $k_{s,t}$ distinct sites. Young's inequality
absorbs the square root of $M_{s,t}$ and proves \eqref{eq:interval}.
The choice $s=0,t=n$ proves \eqref{eq:M}.

To pass from the lattice sum to $U_{D_n}$, first replace $Db(x-y)$
by $(2/\omega_d)(x-y)/|x-y|^d$. The errors sum to
$C\sum_{|x|\leq n}(1+|x-y|)^{-d}\leq CL$.
Replacing each kernel value by its cell integral has the same bound:
the distant-cell error is $O_d((1+|x-y|)^{-d})$, and the kernel
$|v-y|^{1-d}$ is locally integrable near the remaining cells.
Since $|u_x-u_v|\leq C_d/(1+|v|)$ on $C_x$, replacing the direction
costs at most
\begin{equation}\label{eq:direction-error}
C_d\int_{B(0,n+\sqrt d)}
\frac{\dd v}{(1+|v|)|v-y|^{d-1}}\leq C_dL\, .
\end{equation}
For this bound, remove $|v-y|<1$ and use H\"older's inequality with
exponents $d$ and $d/(d-1)$; both norm powers are logarithmic.
Moving $y$ inside its cell costs another $O_d(L)$ by the same
near/far decomposition. Equation~\eqref{eq:approx} follows.
\end{proof}

We will retain the variance at individual targets. Since $B_y\leq n$,
dyadic bounds for the bracket in \eqref{eq:bracket} give, on an event
with the same probability bound,
\begin{equation}\label{eq:localmart}
|\mathcal M_n^y|\leq C(\sqrt{B_yL}+L)
\qquad(y\in\Z^d,\ |y|\leq3n)\, .
\end{equation}
The deterministic replacements in the preceding proof also give
\begin{equation}\label{eq:pointwise}
\ell_n(y)=U_{D_n}(y)+\rho_n(y)-\mathcal M_n^y ,\qquad
|\rho_n(y)|\leq CL\, .
\end{equation}
These formulas apply simultaneously to the same deterministic targets.
For real points $y,z$, we also use the cell modulus
\begin{equation}\label{eq:cellmodulus}
|U_{D_n}(y)-U_{D_n}(z)|\leq CL
\qquad(|y-z|\leq\sqrt d,\ |y|,|z|\leq2n)\, .
\end{equation}

\subsection{Spherical averages}
The radial average removes the angular dependence of the potential.
Let $\sigma_d=d\omega_d$ be the surface area of the unit sphere.
For a bounded measurable set $D\subset\R^d$, write
\begin{equation}\label{eq:Fdef}
f(r)=\frac1{\sigma_d}\int_{S^{d-1}}\ind_D(r\theta)\dd S(\theta) ,
\qquad F(s)=\int_s^\infty f(r)\dd r\, .
\end{equation}
Then $F$ is nonincreasing and absolutely continuous, and
\begin{equation}\label{eq:Fmass}
F(s)=\frac1{\sigma_d}\int_{D\cap\{|v|>s\}}|v|^{1-d}\dd v
\leq\frac{|D|}{\sigma_ds^{d-1}}
\qquad(s>0)\, .
\end{equation}

\begin{lemma}\label{lem:geometry}
Let $d\geq2$ be an integer, let $\eps>0$, and let $D\subset\R^d$
be a bounded measurable set of volume $R$. Let
$U_D(y)=(2\eps/\omega_d)\int_Du_v\cdot(v-y)|v-y|^{-d}\dd v$,
and let $F(s)=\sigma_d^{-1}\int_{D\cap\{|v|>s\}}|v|^{1-d}\dd v$.
Then there is a constant $C_d>0$, depending only on $d$, such that
for every $y,z\in\R^d$ and every $s>0$,
\begin{align}
\|U_D\|_\infty&\leq C_d\eps R^{1/d}\, ,\label{eq:potential-bound}\\
|U_D(y)-U_D(z)|&\leq C_d\eps R^{1/(2d)}|y-z|^{1/2}\, ,\label{eq:holder}\\
\frac1{\sigma_d}\int_{S^{d-1}}U_D(s\theta)\dd S(\theta)
&=2d\eps F(s)\, .\label{eq:newton}
\end{align}
For every $b>0$, the ball potential is
\begin{equation}\label{eq:ballpotential}
U_{B(0,b)}(y)=2d\eps(b-|y|)_+\, .
\end{equation}
\end{lemma}

\begin{proof}
Integrating $|v-y|^{1-d}$ over a ball of volume $R$ centered at $y$ proves
\eqref{eq:potential-bound}. For $K(v)=v/|v|^d$ and
$q=2d/(2d-1)$, scaling gives
\begin{equation}\label{eq:kernel-modulus}
\|K(\cdot-y)-K(\cdot-z)\|_{L^q(\R^d)}=C_d|y-z|^{1/2}\, .
\end{equation}
The norm is finite: the singularities have order $d-1$, and the
difference at infinity has order $d$. H\"older's inequality with
the $L^{2d}$ norm of $\ind_Du$ proves \eqref{eq:holder}.

The spherical mean of $\log|v-s\theta|$ in dimension two is
$\log\max(|v|,s)$. For $d\geq3$, the spherical mean of
$|v-s\theta|^{2-d}$ is $\max(|v|,s)^{2-d}$.
These identities follow from the mean-value property and radial
harmonicity inside and outside the source sphere.
Differentiating with respect to $|v|$, away from $|v|=s$, gives
\begin{equation}\label{eq:kernel-average}
\frac1{\sigma_d}\int_{S^{d-1}}
u_v\cdot\frac{v-s\theta}{|v-s\theta|^d}\dd S(\theta)
=|v|^{1-d}\ind_{\{|v|>s\}}\, .
\end{equation}
Fubini's theorem applies because
$\sup_y\int_D|v-y|^{1-d}\dd v<\infty$.
Integration over $D$ proves \eqref{eq:newton}; continuity covers
every radius. For a ball, $U_D$ is radial, and
$F(s)=(b-s)_+$. This proves \eqref{eq:ballpotential}.
\end{proof}
\section{Bounds on the radius and occupation times}\label{sec:coarse}
We first bound the range and local times at their natural scales.

\begin{proposition}\label{prop:coarse}
Let $d\geq2$ be an integer, let $0<\eps<1/d$, and let $(X_j)_{j\geq0}$ be centrally
excited random walk on $\Z^d$, started at the origin, with first-departure probabilities
$q_x(\pm e_i)=1/(2d)\mp\eps u_x^i/2$ and simple random walk steps
on subsequent departures and at the origin. For every $p>0$, there are constants
$c,C>0$, depending only on $d,\eps,p$, such that for every integer
$n\geq2$, with probability at least $1-Cn^{-p}$,
\begin{equation}\label{eq:coarse}
cN^d\leq R_n\leq CN^d ,\qquad
cN\leq M_n\leq CN ,\qquad cN\leq H_n\leq CN\, ,
\end{equation}
where $N=n^{1/(d+1)}$, $R_n$ counts distinct departure sites,
$M_n$ is the maximum departure local time, and
$H_n=\max_{j\leq n}|X_j|$.
\end{proposition}

\subsection{A radial test}
The discrete Laplacian of the truncated potential is supported near
the truncation level.
Throughout this section, $F$ is the tail in \eqref{eq:Fdef} for $D_n$.
Fix $b_d=\lceil\sqrt d/2\rceil+6$, and let
\begin{equation}\label{eq:shell-def}
M_{\rm sh}(r)=\max_{x:||x|-r|\leq3}\ell_n(x)\, .
\end{equation}
The potential kernel $b$ was defined in Section~\ref{sec:local}.
For $r>0$, write
\begin{equation}\label{eq:radial-test}
\phi_r(x)=(b(x)-h_d(r))_+ ,\qquad
h_d(r)=\begin{cases}
\dfrac2\pi\log r+\kappa,&d=2,\\[3pt]
-\dfrac2{(d-2)\omega_d}r^{2-d},&d\geq3
\end{cases}\, .
\end{equation}

\begin{lemma}\label{lem:radial}
Let $d\geq2$ be an integer, let $0<\eps<1/d$, and let $(X_j)_{j\geq0}$ be centrally
excited random walk started at the origin, with first-departure probabilities
$q_x(\pm e_i)=1/(2d)\mp\eps u_x^i/2$, with simple symmetric
steps on later departures and at the origin.
Let $b_d=\lceil\sqrt d/2\rceil+6$, and let
$F(s)=\sigma_d^{-1}\int_{D_n\cap\{|v|>s\}}|v|^{1-d}\dd v$.
For every $p>0$, there are $r_0=r_0(d)>2b_d$ and $C=C(d,\eps,p)>0$
such that, for every integer $n\geq2$, with probability at least
$1-Cn^{-p}$, simultaneously for
all integers $r_0\leq r\leq n$,
\begin{equation}\label{eq:radial}
\begin{split}
F(r+b_d)\leq{}&CM_{\rm sh}(r)
 [F(r-b_d)-F(r+b_d)]\\
&+C\left[\sqrt{M_nr^{1-d}F(r-b_d)L}+r^{1-d}L\right]\, .
\end{split}
\end{equation}
Here $L=\log(n+2)$ and
$M_{\rm sh}(r)=\max_{x:\,||x|-r|\leq3}\ell_n(x)$.
\end{lemma}

\begin{proof}
The asymptotics \eqref{eq:kernel-asymptotics} imply that, for all
sufficiently large $r$,
\begin{equation}\label{eq:levelsets}
\phi_r(x)=0\quad\hbox{if }|x|\leq r-1 ,\qquad
\phi_r(x)=b(x)-h_d(r)\quad\hbox{if }|x|\geq r+1\, .
\end{equation}
Indeed, the radial leading term increases by order $r^{1-d}$
over a unit interval, while its error is $O(r^{-d})$.
The radial upper and lower envelopes are increasing for large radii.
The remaining finite set of inner sites is covered by $h_2(r)\to\infty$
in the plane, and by $h_d(r)\to0$ and $b(x)<0$ in higher dimensions.

Every increment of $\phi_r$ vanishes from $|x|\leq r-2$.
For $|x|>r+3$, exact harmonicity and \eqref{eq:gradient} give
\begin{equation}\label{eq:radial-drift}
(P-I)\phi_r(x)=0 ,\qquad
u_x\cdot D\phi_r(x)\geq c_d|x|^{1-d}\, .
\end{equation}
On the shell $||x|-r|\leq3$, both $|(P-I)\phi_r(x)|$ and
$|D\phi_r(x)|$ are at most $C_dr^{1-d}$, because the positive-part
map is $1$-Lipschitz.

The Dynkin martingale $W^{(r)}$ for $\phi_r(X_j)$ has increments
bounded by $C_dr^{1-d}$ and bracket
\begin{equation}\label{eq:radialbracket}
\langle W^{(r)}\rangle_n
\leq C_d\sum_{|x|>r-2}\ell_n(x)|x|^{2-2d}
\leq C_dM_nr^{1-d}F(r-b_d)\, .
\end{equation}
For the second inequality, each departure cell in the sum lies outside
$r-b_d$, its weighted integral is comparable to $|x|^{1-d}$, and
$|x|^{2-2d}\leq C_dr^{1-d}|x|^{1-d}$.
The martingale $r^{d-1}W^{(r)}$ has bounded increments and bracket
at most $C_dn$. Dyadic variance bounds, including the bound below one,
and a union bound over integer $r$ give
\begin{equation}\label{eq:radialmart}
|W_n^{(r)}|\leq C\left[
\sqrt{M_nr^{1-d}F(r-b_d)L}+r^{1-d}L\right]\, .
\end{equation}

Since $\phi_r(0)=0$ and $\phi_r(X_n)\geq0$, Dynkin's formula gives
\begin{equation}\label{eq:radial-source}
c_d\eps\sum_{\substack{x\in A_n\\|x|>r+3}}|x|^{1-d}
\leq C_dr^{1-d}\sum_{||x|-r|\leq3}\ell_n(x)+|W_n^{(r)}|\, .
\end{equation}
The first-departure contribution on the shell is absorbed by its local
time. The left side bounds a positive constant times $F(r+b_d)$.
Every shell cell lies between $r-b_d$ and $r+b_d$, so their number
is at most
\begin{equation}\label{eq:shell-count}
\sigma_d(r+b_d)^{d-1}[F(r-b_d)-F(r+b_d)]\, .
\end{equation}
Multiplying by $r^{1-d}$ proves \eqref{eq:radial}.
\end{proof}

\subsection{Reducing the exterior volume}
The spherical average and H\"older continuity bound the largest local
time on a shell in terms of the exterior volume.
Write $s_*=R_n^{1/d}$ and $\delta=Ce_n(M_n)$.
Suppose for now that
\begin{equation}\label{eq:preconditions}
s_*\geq cN ,\qquad M_n\leq Cs_* ,\qquad
\delta\leq C\sqrt{s_*}L\, .
\end{equation}
We verify these estimates before using the argument.
Let
\begin{equation}\label{eq:alpha-beta}
\alpha=\frac1{2d-1} ,\qquad \beta=\frac{d-1}{2d-1}\, .
\end{equation}
For a fixed $B_0\geq2$ and $r\in[s_*,B_0s_*]$, we claim that
\begin{equation}\label{eq:shell}
M_{\rm sh}(r)\leq CB_0^\beta s_*^{1-\alpha}
 [F(r-b_d)+\delta]^\alpha\, ,
\end{equation}
with $C$ independent of $B_0$.
On the event of Lemma~\ref{lem:local}, $U_{D_n}+\delta\geq0$
throughout the approximation region.
Let $h$ be its maximum on a sphere with
$s_*/2\leq t\leq B_0s_*+3$. For fixed $B_0$, these spheres lie
inside $B(0,2n)$ for large $n$.
By \eqref{eq:holder}, the function is at least $h/2$ on a spherical
cap of chord radius $ch^2/s_*$. The constant $c$ can be chosen so
that this radius is at most $t/2$, since $h\leq Cs_*$.
The cap area is at least $c_d(h^2/s_*)^{d-1}$. Hence
\begin{equation}\label{eq:cap-average}
2d\eps F(t)+\delta\geq
\frac{ch^{2d-1}}{s_*^{d-1}t^{d-1}}\, .
\end{equation}
Apply this bound on the sphere through each shell site and use the
monotonicity of $F$. This proves \eqref{eq:shell}.

We next reduce $F$ to the threshold
\begin{equation}\label{eq:floor}
\Lambda=Cs_*^{2-d}L\, .
\end{equation}
Use a grid of spacing $2b_d$, with integer $r$ at each midpoint.
For $r\geq s_*$, the error in \eqref{eq:radial} is at most
\begin{equation}\label{eq:halving-error}
C\sqrt{s_*^{2-d}F(r-b_d)L}+Cs_*^{1-d}L\, .
\end{equation}
If $F(r-b_d)\leq A$ and $A\geq\Lambda$, then the error is at most
$A/4$, after increasing the constant in \eqref{eq:floor}.
Unless $F(r+b_d)\leq A/2$, the decrease is at least $cA/K(A)$,
where
\begin{equation}\label{eq:halving-K}
K(A)=CB_0^\beta s_*^{1-\alpha}(A+\delta)^\alpha+1\, .
\end{equation}
A halving therefore takes at most $CK(A)$ grid steps.
For an initial bound $A_0\leq Cs_*$, the number of dyadic levels
above $\Lambda$ is $O_d(L)$, because $s_*\leq n^{1/d}$.
The sum of their $A^\alpha$ contributions is bounded by
$CA_0^\alpha$. Thus the radial distance required to reach $\Lambda$
is at most
\begin{equation}\label{eq:halving-cost}
CB_0^\beta s_*^{1-\alpha}A_0^\alpha
+C_{B_0}s_*^{1-\alpha/2}L^{1+\alpha}+CL\, .
\end{equation}
The first constant is independent of $B_0$.
The last two terms are $o(s_*)$, uniformly under \eqref{eq:preconditions}.
This calculation is valid whenever the stated distance fits inside
the window $[s_*,B_0s_*]$.

\subsection{Crossings through a small set}
An interval with few distinct departure sites cannot make a long
outward crossing.
The compensated position
\begin{equation}\label{eq:vector-def}
Z_t=X_t+\eps\sum_{j<t}I_ju_{X_j}\, .
\end{equation}
is a vector martingale with bounded increments.
Azuma's inequality in each coordinate and a union bound over deterministic
intervals give, with probability at least $1-Cn^{-p}$,
\begin{equation}\label{eq:vector}
|Z_t-Z_s|\leq C\sqrt{(t-s)L}
\qquad(0\leq s<t\leq n)\, .
\end{equation}
Consequently every projection chosen from the path is also controlled.

Consider an interval $[s,t]$ whose departure sites satisfy
$v\cdot X_j>b>0$ and $|X_j|<r$, where $|v|=1$ and $b/r\geq\kappa>0$.
Suppose that $v\cdot(X_t-X_s)\geq h>0$.
Let $k$ be the number of fresh departures and $m$ the number of
distinct departure sites on the interval.
Equation~\eqref{eq:interval} bounds its duration by
$m(C\eps k^{1/d}+C\lambda_n)$.
Each fresh departure has projected inward mean at least $\kappa\eps$.
By \eqref{eq:vector},
\begin{equation}\label{eq:crossing-before-young}
(h+\kappa\eps k)^2\leq Cm(\eps k^{1/d}+\lambda_n)L\, .
\end{equation}
Young's inequality with powers $2d$ and $2d/(2d-1)$ absorbs half
the $(\kappa\eps k)^2$ term and gives
\begin{equation}\label{eq:crossing}
h^2\leq C_\kappa\bigl(m^\gamma L^\gamma+m\lambda_nL\bigr) ,
\qquad\gamma=\frac{2d}{2d-1}<2\, .
\end{equation}
The event on which this implication is used is simultaneous over
all deterministic intervals and does not depend on the later choice
of $v,b,r$.

\subsection{Proof of the occupation and radius bounds}
We now close the estimates without a prior bound on the outer radius.

\begin{proof}[Proof of Proposition~\ref{prop:coarse}]
Intersect the events in Lemmas~\ref{lem:local} and~\ref{lem:radial}
with the event in \eqref{eq:vector}.
Since $n\leq M_nR_n$, equation~\eqref{eq:M} gives
\begin{equation}\label{eq:sstar-lower}
n\leq C\eps s_*^{d+1}+C\lambda_ns_*^d\, .
\end{equation}
This forces $s_*\geq cN$ for large $n$: if $c$ is sufficiently small,
then the first term for $s_*<cN$ is at most $n/2$ and the second is $o(n)$.
It follows that $\lambda_n=o(s_*)$, $M_n\leq Cs_*$, and
$\delta\leq C\sqrt{s_*}L$. This verifies \eqref{eq:preconditions}.

Start the halving argument at $\lceil s_*\rceil$, where
$F\leq s_*/\sigma_d$.
Choose $B_0\geq8$ so large that the leading term
$CB_0^\beta s_*$ in \eqref{eq:halving-cost} is less than $B_0s_*/4$.
This is possible because $\beta<1$.
For sufficiently large $n$, the remaining terms are also less than
$B_0s_*/4$. The halvings therefore finish before $(B_0-2)s_*$, and
\begin{equation}\label{eq:coarse-tail}
F((B_0-2)s_*)\leq Cs_*^{2-d}L\, .
\end{equation}
All integer radii used lie below $n$ for large $n$, since
$s_*\leq n^{1/d}$ and $B_0$ is fixed.

Suppose for a contradiction that $H_n\geq(B_0+1)s_*$.
Let $\tau$ be the first time the path reaches or exceeds this radius, let
$v=X_\tau/|X_\tau|$, and let $b=(B_0-1)s_*$.
Choose the last entrance into the half-space $v\cdot x>b$ before
$\tau$:
\begin{equation}\label{eq:coarse-entrance}
s=1+\max\{j<\tau:v\cdot X_j\leq b\}\, .
\end{equation}
The set contains zero. At entrance, $v\cdot X_s\leq b+1$.
For $s\leq j<\tau$, the departure site $X_j$ has projection greater
than $b$ and radius less than $(B_0+1)s_*$. The net projected displacement
is at least $s_*$ for large $n$.
Their $m$ distinct cells belong to $D_n$, including when $\tau=n$.
By \eqref{eq:Fmass} and \eqref{eq:coarse-tail},
\begin{equation}\label{eq:coarse-cellcount}
m\leq\sigma_d((B_0+1)s_*+\sqrt d/2)^{d-1}
F(b-\sqrt d/2)\leq C_{B_0}s_*L\, .
\end{equation}
The crossing bound \eqref{eq:crossing} would imply
\begin{equation}\label{eq:coarse-contradiction}
s_*^2\leq C_{B_0}\bigl(s_*^\gamma L^{2\gamma}+s_*L^4\bigr)
=o(s_*^2)\, .
\end{equation}
We conclude that
\begin{equation}\label{eq:HvsR}
H_n\leq CR_n^{1/d}\, .
\end{equation}

It remains to bound $R_n$ in terms of $n$.
Let $S=H_n+\sqrt d$, so that $D_n\subset B(0,S)$.
Integrating the spherical average in \eqref{eq:newton} gives
\begin{equation}\label{eq:massidentity}
\int_{B(0,S)}U_{D_n}(y)\dd y
=2d\eps\sigma_d\int_0^S r^{d-1}F(r)\dd r
=2\eps\int_{D_n}|v|\dd v\, .
\end{equation}
Since $\int\widetilde\ell_n=n$, the approximation error gives
\begin{equation}\label{eq:coarse-masserror}
\left|n-2\eps\int_{D_n}|v|\dd v\right|
\leq\omega_dS^d\delta\leq Cs_*^{d+1/2}L=o(s_*^{d+1})\, .
\end{equation}
Among sets of volume $R_n$, the centered ball minimizes the integral
of $|v|$. Therefore
\begin{equation}\label{eq:radial-packing}
\int_{D_n}|v|\dd v\geq
\frac d{d+1}\omega_d^{-1/d}R_n^{1+1/d}\, .
\end{equation}
Absorb the error in \eqref{eq:coarse-masserror} to obtain $R_n\leq CN^d$.
Equations~\eqref{eq:M} and~\eqref{eq:HvsR} imply $M_n,H_n\leq CN$.
The lower bounds follow from $R_n\geq cN^d$,
$R_n\leq C_d(H_n+1)^d$, and $n\leq M_nR_n$.
Increasing constants covers bounded $n$.
\end{proof}
\section{The inner radius and the local-time profile}\label{sec:contact}
We compare the range with its largest centered ball and then determine
the radius from the total elapsed time.

Fix $p>0$ and intersect the events in Proposition~\ref{prop:coarse},
Lemmas~\ref{lem:local} and~\ref{lem:radial}, and
\eqref{eq:localmart} and~\eqref{eq:vector}. The complement has probability
at most $Cn^{-p}$. All arguments below are deterministic on this event,
apart from the additional concentration estimate for the quadratic
martingale, whose exceptional probability has the same bound.
Choose a deterministic $K>a+1$ large enough that the coarse radius
bound implies $D_n\subset B(0,(K-1)N)$. On $B(0,KN)$,
\eqref{eq:approx} gives
\begin{equation}\label{eq:global}
|\widetilde\ell_n(y)-U_{D_n}(y)|\leq\delta_0 ,\qquad
\delta_0=\begin{cases}C\sqrt N L,&d=2,\\
C(\sqrt{NL}+L),&d\geq3
\end{cases}\, .
\end{equation}


Let $D=D_n$ and let $W=NQ^{1/d}$. Write the centered inradius and excess volume as
\begin{equation}\label{eq:bm}
 b=\inf\{|y|:y\notin D\},\qquad E=D\setminus B(0,b),\qquad m=|E|\, .
\end{equation}
The ball $B(0,b)$ is contained in $D$, up to immaterial cell boundaries.
Choose $y_0$ of norm $b$ that is a limit of points outside $D$.
There is an unoccupied lattice cell $C_z$ whose closure contains $y_0$.
Thus
\begin{equation}\label{eq:contactcell}
 \ell_n(z)=0,\quad |z-y_0|\leq\sqrt d/2,\quad |z|\geq b\, .
\end{equation}
This follows by taking a convergent subsequence from the finitely
many cells adjacent to $y_0$; no regularity of $\partial D$ is assumed.

\subsection{The potential at an unvisited site}
The excess volume gives a positive contribution to the potential at $y_0$.
We claim that
\begin{equation}\label{eq:contactbound}
c\frac{m}{N^{d-1}}\leq U_{D_n}(y_0)
\leq C(\sqrt{B_zL}+L)\, .
\end{equation}
In addition, for every $y\in\R^d$,
\begin{equation}\label{eq:packing}
|U_E(y)|\leq Cm^{1/d}\, .
\end{equation}

If $s=|v|\geq b=|y_0|$, then
\[
 \frac{u_v\cdot(v-y_0)}{|v-y_0|^d}
 =\frac{s^2-v\cdot y_0}{s|v-y_0|^d}
 \geq\frac1{2s}|v-y_0|^{2-d}
 \geq 2^{1-d}s^{1-d}\, .
\]
The first inequality uses
$2(s^2-v\cdot y_0)=|v-y_0|^2+s^2-b^2$; the second uses
$|v-y_0|\leq2s$ and $d\geq2$.
The singular point is immaterial in the integral. Since
$U_{B(0,b)}(y_0)=0$ and $s\leq CN$ on $D$, integration proves
the lower bound in \eqref{eq:contactbound}. The upper bound follows
from \eqref{eq:pointwise}, \eqref{eq:localmart}, \eqref{eq:cellmodulus}, and
$\ell_n(z)=0$; the target is admissible because $|z|=O(N)<3n$. Finally $|u_v|\leq1$, and integrating
$|v-y|^{1-d}$ over a set of volume $m$ is maximized by a centered ball.
Its integral is $O(m^{1/d})$, proving \eqref{eq:packing}.



\subsection{The contact variance in higher dimensions}
Suppose that $d\geq3$. A summable convolution bounds the local variance
near the unvisited site.
Put $V=m^{1/d}$ and $s_b(y)=(b-|y|)_+$. Equations
\eqref{eq:pointwise}, \eqref{eq:localmart}, \eqref{eq:ballpotential}, and
\eqref{eq:packing} imply, at occupied lattice sites,
\[
 \ell_n(y)\leq c_0s_b(y)+C(V+L)+C\sqrt{B_yL}\, .
\]
Let
\[
 k(x)=(1+|x|)^{2-2d}\ind_{\{|x|\leq3n\}},\quad
 k_0=\sum_xk(x),\quad J=\sum_x|x|k(x)\, .
\]
For sufficiently large $n$, \eqref{eq:coarse} ensures that
$B_y=(k*\ell_n)(y)$ at every occupied site and at $z$.
Moreover
\begin{equation}\label{eq:kernelmoments}
 k_0\leq C,\qquad
 J\leq\begin{cases}CL,&d=3,\\C,&d\geq4\end{cases}\, .
\end{equation}
Choose a deterministic $\lambda>0$ so small that
$\lambda k_0<1/2$ for every $n$. Young's inequality gives
\begin{equation}\label{eq:resolvent}
 \ell_n\leq c_0s_b+C(V+L)+\lambda k*\ell_n
 \quad\hbox{on all of }\Z^d\, .
\end{equation}
Outside the occupied set the inequality holds trivially because its
left side is zero and its right side is nonnegative.

Since $s_b$ is $1$-Lipschitz,
\[
 (k^{*j}*s_b)(y)\leq k_0^j s_b(y)+j k_0^{j-1}J\qquad(j\geq1)\, .
\]
Iterate \eqref{eq:resolvent} and sum the two geometric series. The remaining
term tends to zero since the operator norm on $\ell^\infty$ is
$\lambda k_0<1/2$ and $\ell_n$ is bounded.
Using \eqref{eq:kernelmoments} yields
\begin{equation}\label{eq:envelopehigh}
 \ell_n(y)\leq C(s_b(y)+V+L)\quad(y\in\Z^d)\, .
\end{equation}
At the unoccupied center $z$, $s_b(z)=0$, so convolution once more gives
\begin{equation}\label{eq:holebracket}
 B_z\leq C(V+L)\, .
\end{equation}
Indeed $s_b(x)\leq|x-z|$ and hence $(k*s_b)(z)\leq J$.
The contact inequality now becomes
\[
 m\leq CN^{d-1}\big(m^{1/(2d)}\sqrt L+L\big)\, .
\]
If the second term is at least half the left side, then
$m\leq CN^{d-1}L$. Otherwise raise the inequality for
$m^{1-1/(2d)}$ to the power $2d/(2d-1)$. In both cases,
\begin{equation}\label{eq:masshigh}
 m\leq CN^{2d(d-1)/(2d-1)}L^{d/(2d-1)}
       =CN^d Q\, .
\end{equation}
The bound $N^{d-1}L$ is absorbed because $N/L\to\infty$.

\subsection{The planar contact variance}
Suppose that $d=2$. A first bound on the excess volume removes the
logarithmic loss in the variance at the unvisited site.
The coarse global approximation \eqref{eq:global} and the contact
inequality first give
\[
 m\leq CN^{3/2}L\, .
\]
By \eqref{eq:ballpotential} and \eqref{eq:packing}, for every occupied $x$,
\[
 \ell_n(x)\leq c_0(b-|x|)_+
       +CN^{3/4}L^{1/2}+C\sqrt N L\, .
\]
At the unoccupied center $z$, $(b-|x|)_+\leq|x-z|$.
All occupied sites lie within distance $CN$ of $z$, so
\[
 \sum_{|x-z|\leq CN}\frac{(b-|x|)_+}{(1+|x-z|)^2}
 \leq CN,\qquad
 \sum_{|x-z|\leq CN}(1+|x-z|)^{-2}\leq CL\, .
\]
Therefore
\[
 B_z\leq C\big(N+N^{3/4}L^{3/2}+\sqrt N L^2\big)
       \leq CN\, .
\]
for all sufficiently large $n$. Applying
\eqref{eq:contactbound} again proves
\begin{equation}\label{eq:massplanar}
 m\leq CN^{3/2}\sqrt L=CN^2Q\, .
\end{equation}
Combining this with \eqref{eq:global} and \eqref{eq:packing} gives
\begin{equation}\label{eq:envelopeplanar}
 \ell_n(y)\leq c_0(b-|y|)_++CW,
 \qquad W=N^{3/4}L^{1/4}\, .
\end{equation}
Here $\sqrt N L=o(W)$.
In every dimension, \eqref{eq:envelopehigh} or \eqref{eq:envelopeplanar}
therefore implies
\begin{equation}\label{eq:shellW}
 M_{\rm sh}(r)\leq CW\quad\hbox{whenever }r\geq b+3\, .
\end{equation}

\subsection{Determining the radius}
The squared displacement determines the radius through the elapsed time.
Since every step has squared length one and conditional mean
$-\eps I_ju_{X_j}$, the process
\[
\mathcal Q_t=\sum_{j<t}2X_j\cdot(X_{j+1}-X_j+\eps I_ju_{X_j})\, ,
\]
is a martingale, and
\begin{equation}\label{eq:quadratic}
 |X_n|^2=n-2\eps\sum_{x\in A_n}|x|+\mathcal Q_n\, .
\end{equation}
Stop this martingale upon first exiting $B(0,KN)$. Its increments are bounded by $CN$ and its
bracket by $CnN^2$. Freedman's inequality, intersected with the
coarse event, yields
\begin{equation}\label{eq:quadraticerror}
 |\mathcal Q_n|\leq C\big(N\sqrt{nL}+NL\big)\, .
\end{equation}
with another failure probability at most $Cn^{-p}$.
On the coarse event the stopped and original martingales coincide.

The cell replacement satisfies
\[
 \left|\sum_{x\in A_n}|x|-\int_D|v|\dd v\right|
 \leq (\sqrt d/2)|A_n|\leq CN^d\, .
\]
Also
\[
 \int_D|v|\dd v
 =\frac{d\omega_d}{d+1}b^{d+1}+\int_E|v|\dd v,
 \qquad 0\leq\int_E|v|\dd v\leq CNm\leq CN^{d+1}Q\, .
\]
Divide \eqref{eq:quadratic} by $n=N^{d+1}$. Equations
\eqref{eq:quadraticerror}, \eqref{eq:masshigh}, and \eqref{eq:massplanar} give
\[
 \left|1-\frac{c_0\omega_d}{d+1}(b/N)^{d+1}\right|
 \leq C\left(Q+N^{-1}+N^{-(d-1)/2}\sqrt L
                         +N^{-d}L+N^{1-d}\right)
 \leq CQ\, .
\]
For $d=2$, the largest additional term is $N^{-1/2}\sqrt L=Q$.
For $d\geq3$, it is $O(N^{-1}\sqrt L)=O(Q)$.
Since the positive root of
$c_0\omega_d t^{d+1}/(d+1)=1$ is $a$, this proves
\begin{equation}\label{eq:inradius}
 |b/N-a|\leq CQ\, .
\end{equation}
In particular $b\asymp N$. It follows immediately from
$D=B(0,b)\cup E$ up to null sets that \eqref{eq:volume} holds.
The inner inclusion in \eqref{eq:sandwich} follows by increasing its
constant to avoid cell-boundary equality.

Finally, on the approximation region, \eqref{eq:global},
\eqref{eq:ballpotential}, \eqref{eq:packing}, and \eqref{eq:inradius} give
\[
 \left|\widetilde\ell_n(y)-c_0(aN-|y|)_+\right|
 \leq C(\delta_0+NQ+NQ^{1/d})\leq CW\, .
\]
This proves \eqref{eq:profile-rate}; outside a fixed ball containing the range
both the local time and the comparison cone vanish.

\section{The outer fluctuation bound}\label{sec:outer}
We bound the weighted exterior volume and use the crossing estimate
to exclude excursions beyond the claimed radius.
Start at an integer radius just above $b+2b_d+6$. Since $b\asymp N$,
\[
 F(b)\leq CmN^{1-d}\leq CNQ\, .
\]
All subsequent radii remain in a fixed annulus of scale $N$.
Let $\Lambda=CN^{2-d}L$, choosing $C$ sufficiently large.
If a current dyadic upper bound $A$ satisfies $A\geq\Lambda$, then the
error in \eqref{eq:radial}, when $F(r-b_d)\leq A$, is at most $A/4$.
If $F(r+b_d)>A/2$, then \eqref{eq:radial} and \eqref{eq:shellW} therefore force
\[
 F(r-b_d)-F(r+b_d)\geq \frac{cA}{W+1}\, .
\]
On the grid of spacing $2b_d$, a halving takes at most $C(W+1)$
steps. There are at most $CL$ levels between $CNQ$ and
$\Lambda$. If the initial bound is already below $\Lambda$, then no
halving is needed. Thus, for a deterministic constant $K_0$,
\begin{equation}\label{eq:tailend}
 F(b+K_0WL)\leq\Lambda\, .
\end{equation}
Rounding errors are absorbed because $W\to\infty$.
Since $WL=o(N)$, all required radii fit in that fixed annulus.

Set $h=K_1WL$, with $K_1>2K_0$, and suppose that the path
reaches radius $r=b+3h$ by time $n$. Let
\[
 \tau=\min\{t\leq n:|X_t|\geq r\},\quad
 v=X_\tau/|X_\tau|,\quad \beta=b+2h,
 \quad s=1+\max\{j<\tau:v\cdot X_j\leq\beta\}\, .
\]
The set defining $s$ contains zero. All departures from $s$ through
$\tau-1$ have projection greater than $\beta$ and radius less than
$r$, while their net projected displacement is at least $h-1$.
The ratio $\beta/r$ is bounded below by a fixed positive constant.
The distinct departure cells in this interval lie in $D_n$ even
if $\tau=n$. By \eqref{eq:tailend}, their number $m_{\rm cap}$ satisfies
\[
 m_{\rm cap}\leq CN^{d-1}F(\beta-\sqrt d/2)
              \leq CNL\, .
\]
Consequently \eqref{eq:crossing} would imply
\begin{equation}\label{eq:outer-contradiction}
 (h-1)^2\leq C
 \begin{cases}
 N^{4/3}L^{8/3}+NL^4,&d=2,\\
 N^{2d/(2d-1)}L^{4d/(2d-1)}+NL^3,&d\geq3
 \end{cases}\, .
\end{equation}
The right side is $o(h^2)$: for $d=2$,
$h^2\asymp N^{3/2}L^{5/2}$; for $d\geq3$,
$h^2\asymp N^{(4d-4)/(2d-1)}L^{4d/(2d-1)}$ and
$4d-4>2d$. This contradiction proves the outer inclusion in
\eqref{eq:sandwich}, using \eqref{eq:inradius} and $NQ=O(WL)$.

The lattice mesh contributes at most $\sqrt d/N$ to Hausdorff distance,
which is absorbed in \eqref{eq:hausdorff}.



\begin{proof}[Proof of Theorems~\ref{thm:shape} and~\ref{thm:fluctuations}]
The preceding estimates prove the inequalities in
Theorem~\ref{thm:fluctuations} at each deterministic $n$, with total
failure probability at most $Cn^{-p}$. Fix $p>1$. Borel--Cantelli
implies that these inequalities hold for every sufficiently large
integer $n$, almost surely. Since $Q\to0$ and $Q^{1/d}L\to0$,
the ball inclusions and profile convergence in Theorem~\ref{thm:shape}
follow. The volume estimate gives $|V_n|/N^d\to\omega_da^d$ because
$|V_n\setminus A_n|\leq1$. For each fixed lattice site $x$,
the uniform profile estimate and $|x|/N\to0$ give
$\ell_n(x)/N\to2d\eps a>0$. The lattice is countable, so the
walk visits every vertex infinitely often on the same probability-one event.
\end{proof}

\bibliographystyle{plainnat}
\bibliography{references/references}
\end{document}
